\section{Elementary Koenigs functions at $r=2$ and $r=4$}
\label{a:app:integrable}
These cases lie outside subcritical branching ($r=2p<1$) but complete the classical picture of integrable logistic linearisations. The derivations follow the closed-form note integrated into this document.

\subsection{Case $r=2$ (logarithmic)}

For $r=2$, $f_2(x)=2x(1-x)$. The fixed point $0$ has multiplier $\lambda=2$ (repelling); the linearisation is of Poincar\'e type rather than an attracting-Koenigs function, but an elementary closed form still exists.

The affine change $u=1-2x$ conjugates the dynamics to $u\mapsto u^2$. Linearising $u\mapsto u^2$ by the logarithm and returning to $x$ yields the candidate
\begin{equation}
\label{a:eq:phi2}
\psi_2(x)=-\frac12\ln(1-2x),
\end{equation}
normalised so that $\psi_2(0)=0$ and $\psi_2'(0)=1$.

\textbf{Verification.}
\begin{align*}
\psi_2\bigl(f_2(x)\bigr)
&=
-\frac12\ln\bigl(1-2\cdot 2x(1-x)\bigr)
=
-\frac12\ln\bigl((1-2x)^2\bigr)
=
-\ln(1-2x)
=
2\psi_2(x).
\end{align*}

\subsection{Case $r=4$ (inverse trigonometric)}

For $r=4$, $f_4(x)=4x(1-x)$. The fixed point $0$ has multiplier $f_4'(0)=4$ (repelling); as with $r=2$, this is a Poincar\'e-type linearisation, not an attracting Koenigs function. With $x=\sin^2\theta$ one finds $f_4(x)=\sin^2(2\theta)$, so the angle doubles. The candidate
\begin{equation}
\label{a:eq:phi4}
\psi_4(x)=\bigl(\arcsin\sqrt{x}\bigr)^2
\end{equation}
satisfies $\psi_4(f_4(x))=4\psi_4(x)$ on a suitable real domain, and $\psi_4(x)\sim x$ as $x\to0$, so $\psi_4'(0)=1$.

\textbf{Verification.}
Using $\arcsin(2u\sqrt{1-u^2})=2\arcsin u$ with $u=\sqrt{x}$ (on an appropriate interval),
\begin{align*}
\psi_4\bigl(f_4(x)\bigr)
&=
\bigl(\arcsin\sqrt{4x(1-x)}\bigr)^2
=
\bigl(2\arcsin\sqrt{x}\bigr)^2
=
4\psi_4(x).
\end{align*}

\subsection{Affine conjugacy to $z^2+c$}

\begin{lemma}
\label{a:lem:affine}
The logistic map $f_r(x)=rx(1-x)$ is affinely conjugate to $P_c(z)=z^2+c$ with $c=(2r-r^2)/4$.
\end{lemma}

\begin{proof}[Proof sketch]
Seek $h(x)=ax+b$ with $h\circ f_r\circ h^{-1}=P_c$. The substitution $x=-z/r+\tfrac12$ (for $r\neq0$) yields a monic quadratic; matching coefficients gives $c=(2r-r^2)/4$.
\end{proof}

Under this conjugacy, $r=2\Leftrightarrow c=0$ (monomial $z^2$) and $r=4\Leftrightarrow c=-2$ (Chebyshev type; the value $c=-2$ is also reached at $r=-2$). Neither exceptional value lies in the image of $r\in(0,1)$, which is $c\in\bigl(0,\tfrac14\bigr)$. As explained in \cref{a:rem:empty}, this parameter count is a consistency check rather than the proof: the decisive point is that both exceptional cases sit over \emph{repelling} fixed points, which the attracting window excludes structurally.

\subsection{The worked examples inside the Becker--Bergweiler classification}

The precise classification is \cref{a:thm:BB-main} in the main text: a differentially algebraic Schr\"oder function exists only over a repelling fixed point, and is then a M\"obius transformation of $\exp(\alpha t^{\rho})$, of $\cos(\alpha t^{\rho}+\beta)$, or of Weierstrass-$\wp$ material, with the underlying map M\"obius-conjugate to $z^{\pm d}$, $\pm T_d$, or a Latt\`es map. The two derivations above realise the first two families exactly.

\emph{Case $r=2$.} Inverting \eqref{a:eq:phi2}: $t=-\tfrac12\ln(1-2x)$ gives
\[
\varphi_2(t)=\psi_2^{-1}(t)=\frac{1-e^{-2t}}{2},
\]
an affine (hence M\"obius) image of $\exp(-2t)$ --- family $\exp(\alpha t^{\rho})$ with $\rho=1$, $\alpha=-2$. Correspondingly the substitution $u=1-2x$ conjugates $f_2$ to the monomial $u\mapsto u^2$, and the linearised fixed point has multiplier $2$ (repelling), as \cref{a:thm:BB-main}(i) requires.

\emph{Case $r=4$.} Inverting \eqref{a:eq:phi4}: $t=(\arcsin\sqrt{x})^2$ gives
\[
\varphi_4(t)=\psi_4^{-1}(t)=\sin^2\!\sqrt{t}=\frac{1-\cos\bigl(2\sqrt{t}\bigr)}{2},
\]
a M\"obius image of $\cos(\alpha t^{\rho}+\beta)$ with $\rho=\tfrac12$, $\alpha=2$, $\beta=0$ --- the fractional exponent is admissible precisely because $\cos(2\sqrt{t})$ is an entire function of $t$. Correspondingly $u=1-2x$ conjugates $f_4$ to $T_2(u)=2u^2-1$, the second Chebyshev polynomial, and the multiplier at the linearised fixed point is $4$ (repelling).

Both examples therefore sit exactly where \cref{a:thm:BB-main} says all differentially algebraic Schr\"oder functions must sit: over repelling fixed points, in the exponential and cosine families. For $r\in(0,1)$ the fixed point is attracting, no member of the classification is available, and \cref{a:thm:main-ht} gives hypertranscendence of $\psi_r$ --- which underpins the main-text claims for $A(p)=2\psi_{2p}(\tfrac12)$.
